\documentclass[11pt]{article}

\usepackage[utf8]{inputenc}
\usepackage[T1]{fontenc}
\usepackage[margin=2.4cm]{geometry}
\usepackage{xcolor}
\usepackage[skins,breakable]{tcolorbox}
\usepackage{enumitem}
\usepackage{titlesec}
\usepackage{hyperref}
\usepackage{longtable}
\usepackage{booktabs}

% ---------------------------------------------------------------
% Same palette as the slide deck, for visual continuity
% ---------------------------------------------------------------
\definecolor{sovDark}{RGB}{18,32,51}
\definecolor{sovGold}{RGB}{214,164,64}
\definecolor{sovRed}{RGB}{163,44,58}
\definecolor{sovGrey}{RGB}{92,102,117}
\definecolor{sovPaper}{RGB}{250,248,244}
\definecolor{sovGreen}{RGB}{58,110,89}

\hypersetup{
  colorlinks=true,linkcolor=sovDark,urlcolor=sovDark,citecolor=sovDark,
  pdftitle={Speaker Script -- Digital Sovereignty for Kashmir},
  bookmarksopen=true
}

\titleformat{\section}{\normalfont\Large\bfseries\color{sovDark}}{\thesection}{0.6em}{}
\titleformat{\subsection}[block]{\normalfont\large\bfseries\color{sovDark}}{}{0em}{}
  [{\color{sovGold}\titlerule[1.2pt]}]
\titlespacing*{\subsection}{0pt}{18pt}{6pt}

\newtcolorbox{sayit}{
  enhanced,breakable,colback=sovDark,colframe=sovGold,coltext=sovPaper,
  boxrule=1pt,arc=2pt,left=8pt,right=8pt,top=6pt,bottom=6pt,
  fonttitle=\bfseries,title={SAY THIS},before skip=6pt,after skip=10pt
}

\newcommand{\onscreen}[1]{\noindent\textbf{\color{sovGrey}On screen:} #1\par\vspace{4pt}}
\newcommand{\timing}[2]{\hfill{\normalsize\color{sovGrey}[\,#1\ \textbar\ running total #2\,]}}

\newenvironment{terms}
  {\noindent\textbf{\color{sovGrey}Terms to know:}
   \begin{description}[leftmargin=0pt,style=sameline,labelindent=0pt,itemsep=2pt,topsep=2pt,font=\normalfont\bfseries\color{sovDark}]}
  {\end{description}\vspace{2pt}}

\title{\color{sovDark}\textbf{Report on Script}\\[4pt]\Large Digital Sovereignty for Kashmir}

\begin{document}
\maketitle

\clearpage
\section{Section 0 --- Why This, Why Now}

\subsection*{Slide 1 --- Title \timing{0:25}{0:25}}
\onscreen{Title, subtitle, your name, ``Open Library Kashmir (OLK)'', July 2026.}
\begin{sayit}
Good evening, everyone. Thank you for coming. Tonight I want to talk about something
we've touched on before in smaller conversations --- building our own digital
infrastructure, right here, so that our data, our platforms, and our future don't have
to live on someone else's servers, under someone else's laws. I'll walk through why this
matters, what it actually takes technically, and end with a real, fundable first step we
could start this year.
\end{sayit}

\subsection*{Slide 2 --- Tonight's Roadmap \timing{0:45}{1:10}}
\onscreen{The table of contents: the six sections plus ``The Hope''.}
\begin{sayit}
Here's the shape of tonight's talk. We'll start with the ``why'' --- the motivation.
Then we'll get practical: the internet and electricity we'd need, the physical room the
servers live in, starting with a VPS as our on-ramp, the hardware for scaling to a
hundred thousand users, and the software and skills we'd need to run it. And I'll close
with a proposal --- a modest, real pilot we could begin, not a someday dream.
\end{sayit}

\subsection*{Slide 3 --- Divider: Why This, Why Now \timing{0:10}{1:20}}
\onscreen{Section title only.}
\begin{sayit}
So --- why this, why now.
\end{sayit}

\subsection*{Slide 4 --- The cost of renting our own knowledge \timing{1:35}{2:55}}
\onscreen{Bar chart of global cloud market share (AWS 30\%, Azure 24\%, Google 11\%,
others 35\%); gold takeaway box on ``whose infrastructure, whose law''.}
\begin{terms}
  \item[CLOUD Act] A 2018 US law: any company headquartered in the US must hand
    government-requested data to US authorities, even if the server holding it is
    physically in another country.
  \item[Hyperscaler] Industry term for a cloud provider operating at massive, global
    scale (AWS, Azure, Google Cloud, Alibaba Cloud, etc.).
\end{terms}
\begin{sayit}
Let's start with an uncomfortable fact. Roughly two-thirds of the world's cloud
infrastructure spend sits with three American companies --- Amazon, Microsoft, and
Google. And every one of them is bound by something called the CLOUD Act --- a US law
that says a US company must hand over data on request, no matter which country the
actual server sits in. So when we say our data is ``hosted in India'' on an American
platform, that's a comforting label --- it is not sovereignty. The question was never
whether we depend on cloud infrastructure. We already do, every day. The real question
is: whose infrastructure, under whose law, and who profits from it.
\end{sayit}

\subsection*{Slide 5 --- OLK today: a working example, and its ceiling \timing{1:20}{4:15}}
\onscreen{Two columns: ``What we already built'' vs.\ ``The ceiling''.}
\begin{terms}
  \item[Next.js / React] The open-source web framework OLK's interface is built with.
  \item[PostgreSQL] The open-source database engine that stores every book, user, and
    request.
  \item[Supabase] An open-source platform that bundles Postgres with authentication,
    file storage, and an API layer --- what OLK currently runs on.
  \item[\texttt{ap-south-1}] Amazon's internal code for its Mumbai data center region.
\end{terms}
\begin{sayit}
Now, here's where we already stand. Open Library Kashmir --- OLK --- is real. People are
donating books, lending them, requesting them, right now. It's built entirely on
open-source technology: Next.js and React for the interface, PostgreSQL as the
database, all wired together through a platform called Supabase. Every layer of that
software, we control. But here's the ceiling: today it physically runs inside Amazon's
data center in Mumbai. Being geographically close isn't the same as being legally ours
--- it's still Amazon's rack, Amazon's power, Amazon's law. The good news is that
because everything is open-source, we don't need to rewrite a single line to change
that. We just need to redeploy it onto infrastructure we actually own.
\end{sayit}

\subsection*{Slide 6 --- Beyond one app: why this could matter for millions \timing{1:05}{5:20}}
\onscreen{Hub-and-spoke diagram: Open Source Infra at the centre, spokes to libraries,
cooperatives, health records, schools.}
\begin{sayit}
And OLK is not the destination --- it's the first baby step. The exact same pattern ---
a self-hosted database, authentication, storage, a web front end --- is the blueprint
for every community institution that doesn't yet have a digital backbone of its own:
farmer and artisan cooperatives, community health records, local school archives. If we
get this right once, we've built a template the whole region can reuse.
\end{sayit}

\section{Section I --- Internet \& Electricity: The Foundation}

\subsection*{Slide 7 --- Divider: Internet \& Electricity \timing{0:10}{5:30}}
\begin{sayit}
Let's talk foundations --- the internet and the electricity this would actually run on.
\end{sayit}

\subsection*{Slide 8 --- How much internet do we actually need? \timing{1:35}{7:05}}
\onscreen{Bandwidth math worked out step by step; takeaway on CDN/edge caching.}
\begin{terms}
  \item[Gbps / Mbps] Gigabits and megabits per second --- units of network speed.
    1 Gbps = 1000 Mbps.
  \item[Concurrent users] People actively using the app \emph{at the same instant},
    as opposed to total registered users.
  \item[CDN / edge cache] A ``Content Delivery Network'' --- copies of your content
    stored on servers closer to users, so they don't have to reach all the way back to
    the origin server for every request.
\end{terms}
\begin{sayit}
Let's do some real math. Say we reach a hundred thousand registered users. At any given
moment, maybe eight percent are actively online --- that's eight thousand concurrent
users. If each one is pulling around thirty kilobytes a second --- normal for a
text-and-image app like ours --- that's roughly two hundred forty megabytes a second, or
about two gigabits per second, sustained. Add burst headroom for peak load, and we're
looking at eight to ten gigabits per second. A typical home or business fiber line here
in Srinagar --- a hundred megabits to a gigabit --- is an order of magnitude short of
that. We'd need datacenter-grade, redundant links --- multiple carriers, multiple
physical paths. And given Kashmir's history with internet shutdowns, a single line isn't
just slow, it's a single point of total failure. Our design principle: keep a cache
outside the valley --- in Hong Kong or Singapore --- so people can still read even if a
local shutdown stops new writes.
\end{sayit}

\subsection*{Slide 9 --- Electricity: the honest picture \timing{1:20}{8:25}}
\onscreen{Two columns: where J\&K's grid stands vs.\ what that means for our design.}
\begin{terms}
  \item[RDSS] Revamped Distribution Sector Scheme --- the central government's ongoing
    power-sector reform program that J\&K is implementing.
  \item[KPDCL] Kashmir Power Distribution Corporation Limited --- the utility running
    the local grid.
  \item[Online double-conversion UPS] A battery backup that continuously converts
    incoming power to DC and back to AC, so it can take over with \emph{zero}
    switchover delay when grid power drops --- unlike cheaper ``standby'' UPS units.
  \item[Auto-transfer switch] Hardware that automatically switches a building's power
    source from grid to generator (and back) without a human flipping anything.
\end{terms}
\begin{sayit}
On power: the grid here is genuinely improving. Billing efficiency has gone from
fifty-six percent to seventy-seven percent under the RDSS reform program. But KPDCL
still runs scheduled maintenance outages regularly, and unscheduled ones remain common,
especially in winter. So our design has to treat grid power as supplementary, not
primary, from day one. That means a rack-level online double-conversion UPS --- a
battery system that instantly and seamlessly takes over the moment grid power drops ---
sized for at least thirty to sixty minutes of runtime. Behind that, a diesel generator
with an automatic transfer switch and enough fuel for a couple of days. Solar and
battery can supplement over time, but never as the sole source --- Kashmir's winter sun
just isn't reliable enough for that.
\end{sayit}

\section{Section II --- Physical Environment}

\subsection*{Slide 10 --- Divider: Physical Environment \timing{0:10}{8:35}}
\begin{sayit}
Now --- where do these machines actually live?
\end{sayit}

\subsection*{Slide 11 --- The room the servers live in \timing{1:20}{9:55}}
\onscreen{Cooling / safety on the left, physical security and a takeaway on the right.}
\begin{terms}
  \item[Free-air / economizer cooling] Using outside air directly to cool a server
    room instead of running mechanical air-conditioning --- works whenever the outside
    air is already cool and dry enough.
  \item[FM-200 / Novec] Clean-agent fire suppression gases that put out fires without
    leaving behind water or residue that would destroy electronics.
  \item[Positive-pressure filtration] Keeping air pressure inside the room slightly
    higher than outside, so filtered air leaks \emph{out} through any gap rather than
    unfiltered dust leaking \emph{in}.
\end{terms}
\begin{sayit}
Kashmir's climate is actually an asset here. Cool air most of the year means we can use
free-air cooling --- essentially just letting outside air do the work --- and only need
real air-conditioning during summer peaks and for humidity control. We'd seal the rack
and use positive-pressure filtration to keep valley dust and pollen out. For safety,
clean-agent fire suppression --- never water sprinklers over live electronics --- plus
smoke and heat detectors wired to automatic shutdown. And physically: one
access-controlled room, logged entry, CCTV, locking cabinets. For our pilot, this is
genuinely one well-prepared room with a single rack --- not a purpose-built data center.
That's a later, bigger project.
\end{sayit}

\section{Section III --- Starting with VPS: Our On-Ramp}

\subsection*{Slide 12 --- Divider: Starting with VPS \timing{0:10}{10:05}}
\begin{sayit}
So how do we actually start? With a VPS --- our on-ramp.
\end{sayit}

\subsection*{Slide 13 --- Why VPS first: the staged path \timing{1:05}{11:10}}
\onscreen{Three-stage diagram: VPS $\rightarrow$ Colocation $\rightarrow$ Owned facility.}
\begin{terms}
  \item[VPS] Virtual Private Server --- a slice of a shared physical machine, rented
    by the hour or month, that behaves like your own private computer.
  \item[Colocation (``colo'')] Renting rack space, power, and network connectivity
    inside someone else's data center, while the physical server hardware inside it is
    yours.
  \item[Capex] Capital expenditure --- upfront money spent buying physical assets
    (as opposed to ongoing rental/subscription costs).
\end{terms}
\begin{sayit}
Think of this as three stages. Stage one: a VPS --- a virtual private server --- rented
from a Chinese cloud provider. Near-zero upfront cost, and it's where we learn real
operations: Linux, backups, monitoring. Stage two: colocation --- we rent rack space,
power, and network in someone else's facility, but the hardware inside it is ours.
That's exactly where our pilot lands. Stage three: our own facility, full physical
control, based here in Kashmir, serving the region. Every one of these stages runs the
exact same open-source software. Moving between them is a migration, not a rewrite.
We're not choosing between a VPS and our own data center --- we're choosing the order in
which we build confidence and capability.
\end{sayit}

\subsection*{Slide 14 --- The Chinese cloud landscape \timing{1:05}{12:15}}
\onscreen{Four bullets on why Chinese providers specifically.}
\begin{sayit}
Why Chinese cloud providers specifically? Alibaba Cloud, Tencent Cloud, and Huawei Cloud
are three of the largest cloud operators in the world by capacity --- mature, globally
reachable, and entirely outside the reach of that US CLOUD Act we talked about earlier.
All three run international sites in Hong Kong and Singapore built specifically for
customers outside mainland China. And Huawei Cloud in particular runs on the same
processor and operating system family --- Kunpeng and openEuler --- that we'd eventually
deploy on our own hardware. One skill set covers both stages. And to be direct about it:
everything from this point forward in the talk is Chinese-made. No AWS, no Azure, no
Google.
\end{sayit}

\subsection*{Slide 15 --- Alibaba vs.\ Tencent vs.\ Huawei Cloud \timing{1:20}{13:35}}
\onscreen{Comparison table across four rows; gold recommendation box.}
\begin{terms}
  \item[ECS] Elastic Compute Service --- Alibaba's (and Huawei's) name for a rentable
    virtual server.
  \item[CVM] Cloud Virtual Machine --- Tencent's name for the same idea.
  \item[Lighthouse] Tencent's simplified, all-in-one VPS product aimed at small
    projects --- easier to set up than a full CVM.
  \item[Latency] The delay for a packet of data to travel from Srinagar to the server
    and back, measured in milliseconds (ms).
\end{terms}
\begin{sayit}
Comparing the three: Alibaba Cloud has the broadest footprint outside China and the most
mature international billing. Tencent Cloud has a product called Lighthouse ---
basically a simplified, all-in-one VPS --- perfect for a first pilot. Huawei Cloud's
advantage is architectural continuity: Kunpeng and openEuler, the same stack as our
future hardware. All three have a nearby region in Hong Kong or Singapore, with latency
to Srinagar around forty to eighty milliseconds --- perfectly usable. My recommendation:
start on Tencent Lighthouse or Alibaba's ECS for the cheapest, fastest first pilot, and
plan our eventual on-premises hardware around Huawei's ecosystem for continuity.
\end{sayit}

\subsection*{Slide 16 --- Getting a Chinese VPS as a foreign user \timing{1:05}{14:40}}
\onscreen{Five numbered practical steps; gold takeaway box.}
\begin{terms}
  \item[ICP license] A mainland Chinese government filing required for any website
    physically hosted on mainland Chinese servers --- avoided entirely by choosing a
    Hong Kong or Singapore region instead.
  \item[Pay-as-you-go vs.\ reserved instance] Paying by actual usage/hour, versus
    committing upfront to a fixed term (e.g.\ one year) for a discount.
\end{terms}
\begin{sayit}
Practically, here's how we'd actually get one of these as outsiders. Register on the
international site --- not the mainland Chinese site. Choose a Hong Kong or Singapore
region --- that alone avoids China's mainland ICP licensing requirement, which only
applies to websites physically hosted inside mainland China. Identity verification is
just a passport and organization documents --- we'd register OLK formally to do this
cleanly. Payment works with an ordinary international credit card. And critically: start
pay-as-you-go on a small instance before committing to any annual plan, so we validate
our actual workload first. None of this requires setting up a China-based business. It
requires an afternoon of paperwork.
\end{sayit}

\section{Section IV --- Hardware for Full Scale: 100,000 Users}

\subsection*{Slide 17 --- Divider: Hardware for Full Scale \timing{0:10}{14:50}}
\begin{sayit}
Now let's talk about the hardware we'd need to actually scale to a hundred thousand
users.
\end{sayit}

\subsection*{Slide 18 --- Storage 101: what we need, how much \timing{1:20}{16:10}}
\onscreen{SSD vs.\ NVMe on the left, sizing math for 100k users on the right.}
\begin{terms}
  \item[SATA SSD] A solid-state drive connected over the older SATA interface,
    capped around 550 MB/s.
  \item[NVMe] A storage protocol that lets SSDs talk directly over the same
    high-speed PCIe lanes as the processor --- much faster than SATA.
  \item[IOPS] Input/Output Operations Per Second --- how many individual read or
    write operations a drive can handle each second; more important than raw speed for
    databases doing lots of small operations.
  \item[WAL] Write-Ahead Log --- PostgreSQL's record of every change, written before
    it's applied, used to guarantee no data is lost in a crash.
\end{terms}
\begin{sayit}
Two storage technologies matter here. SATA SSDs are fast and cheap but capped around
five hundred fifty megabytes a second by their interface. NVMe drives talk directly over
the same high-speed lanes as the processor --- multiple gigabytes per second, and far
more importantly, far higher IOPS, which just means how many small read-and-write
operations it can handle per second. That matters enormously for a database like ours
doing thousands of tiny random reads and writes. For sizing: at a hundred thousand
users, with about five megabytes of data each, plus book cover images, plus database
overhead like indexes and write-ahead logs and backups, we land on a design target of
two to five terabytes of fast NVMe storage, with a separate cheaper tier for older,
archived data.
\end{sayit}

\subsection*{Slide 19 --- RAID: redundancy is not optional \timing{1:20}{17:30}}
\onscreen{Table of RAID 0/1/5/6/10 with fault tolerance and usable capacity; recommendation
box.}
\begin{terms}
  \item[RAID] Redundant Array of Independent Disks --- combining multiple physical
    drives so data survives one (or more) drive failures, and/or goes faster.
  \item[Striping] Splitting data across multiple drives so reads/writes happen in
    parallel --- faster, but no protection on its own.
  \item[Mirroring] Writing the same data to two drives simultaneously --- if one dies,
    the other has an exact copy.
  \item[Parity] Extra, mathematically-derived data stored alongside the real data,
    which lets the array reconstruct a lost drive's contents.
\end{terms}
\begin{sayit}
RAID is how we protect against a hard truth: drives fail. RAID zero just spreads data
across disks for speed, with zero protection --- if one drive dies, everything's gone.
RAID one mirrors two drives --- simple, but you lose half your capacity. RAID five and
six add parity data so the array survives one or two drive failures respectively. RAID
ten mirrors pairs of drives and then stripes across those pairs --- it gives you the best
of both speed and redundancy, at the cost of losing half your raw capacity. Our pick:
RAID ten for the live database --- it's the system of record, worth the extra cost.
RAID six is fine for the bulk image and backup archive, where raw capacity matters more
than speed.
\end{sayit}

\subsection*{Slide 20 --- ECC RAM: why a bit-flip matters \timing{1:05}{18:35}}
\onscreen{Why ECC matters, then sizing guidance for 10--20k and 100k scale.}
\begin{terms}
  \item[ECC RAM] Error-Correcting Code memory --- detects and automatically fixes
    single-bit errors as they happen, instead of silently corrupting data.
  \item[DIMM] Dual In-line Memory Module --- the physical memory stick that plugs
    into a server.
  \item[\texttt{shared\_buffers}] A PostgreSQL setting controlling how much RAM the
    database reserves for its own working cache.
\end{terms}
\begin{sayit}
Here's something people don't think about: ordinary consumer memory has zero protection
against random bit-flips --- caused by things like cosmic rays or electrical noise.
Rare, yes, but real, and it gets worse the more memory chips you run and the longer they
run. ECC --- error-correcting code --- RAM detects and silently fixes single-bit errors
in real time. For a system of record --- who owns which book, who borrowed what --- a
silently corrupted database row is far worse than a slow page load. This is not where we
cut costs. For sizing, a hundred twenty-eight gigabytes gives PostgreSQL comfortable
room for our ten-to-twenty-thousand-user pilot. Full hundred-thousand-user scale would
mean two hundred fifty-six to five hundred twelve gigabytes, or splitting the load
across multiple machines instead.
\end{sayit}

\subsection*{Slide 21 --- The Chinese hardware stack, end to end \timing{1:20}{19:55}}
\onscreen{Table: CPU, chassis, storage, OS, networking, all Chinese vendors; market-share
takeaway.}
\begin{terms}
  \item[Kunpeng] Huawei's family of ARM-based server processors.
  \item[ARM64] A processor architecture (as opposed to the more familiar x86 used by
    Intel/AMD) --- power-efficient, and what Kunpeng chips use.
  \item[TaiShan / Inspur / XFusion] Chinese server manufacturers that build rack
    servers certified to run Kunpeng processors and openEuler.
  \item[YMTC / ZhiTai] Yangtze Memory Technologies Corp --- China's domestic NAND
    flash memory manufacturer; ZhiTai is its consumer/enterprise SSD brand.
  \item[openEuler] Huawei's open-source Linux distribution, built natively for ARM64.
  \item[GbE] Gigabit Ethernet --- wired networking speed, e.g.\ ``10 GbE'' = 10
    gigabits per second.
\end{terms}
\begin{sayit}
So what does an all-Chinese hardware stack actually look like? On the processor side:
Huawei's Kunpeng chips --- ARM-based, with their newest generation shipping up to a
hundred ninety-two cores per chip. Server chassis from Huawei's TaiShan line, or from
Inspur and XFusion --- all of them ship servers certified for Kunpeng and openEuler.
Storage from YMTC, under their ZhiTai brand --- that's China's homegrown 3D NAND flash
memory manufacturer, with their newest drives hitting ten to fifteen gigabytes a second.
And the operating system: openEuler, Huawei's open-source Linux, which already holds
over half the market for newly deployed server operating systems in China. This isn't a
fringe, unproven stack --- Kunpeng-based servers already hold over twenty percent of
China's server market. This is mainstream infrastructure we can simply choose to use.
\end{sayit}

\subsection*{Slide 22 --- A concrete Bill of Materials \timing{1:05}{21:00}}
\onscreen{Parts list table for a 100k-scale node.}
\begin{terms}
  \item[BOM] Bill of Materials --- the itemised parts list for building something.
  \item[PSU] Power Supply Unit.
  \item[RDIMM] Registered DIMM --- an enterprise-grade memory module (more reliable,
    supports larger capacities than consumer memory).
  \item[PDU] Power Distribution Unit --- the rack-mounted power strip that feeds all
    the servers in a rack.
  \item[2U] A unit of rack server height (``2U'' = 2 rack units, about 3.5 inches).
\end{terms}
\begin{sayit}
Putting that into an actual parts list for a hundred-thousand-user-scale machine: a
two-U rack server chassis from TaiShan or XFusion with dual power supplies, one or two
Kunpeng processors, two hundred fifty-six gigabytes of ECC memory, a small mirrored boot
drive, and the real workhorse --- six NVMe drives in RAID ten for the database, giving
us five to eight terabytes of usable, protected space. Dual ten-or-twenty-five-gigabit
networking, openEuler as the operating system, and proper rack-level UPS support. Exact
pricing moves with the currency exchange rate, so treat this as a specification to take
to a reseller for a quote, not a fixed price.
\end{sayit}

\section{Section V --- Software \& the Skills We Need}

\subsection*{Slide 23 --- Divider: Software \& Skills \timing{0:10}{21:10}}
\begin{sayit}
Now --- the software running on top of all that hardware, and the skills our team needs.
\end{sayit}

\subsection*{Slide 24 --- The software stack: what changes \timing{1:20}{22:30}}
\onscreen{``Stays the same'' vs.\ ``New, Chinese, and worth knowing''.}
\begin{terms}
  \item[openGauss] Huawei's own open-source fork of PostgreSQL, built for
    distributed, petabyte-scale workloads.
  \item[TiDB] An open-source distributed SQL database from PingCAP, wire-compatible
    with MySQL, that scales horizontally across many machines.
  \item[OceanBase] Ant Group's distributed database, built for financial-grade
    transaction processing at extreme concurrency.
  \item[OLTP] Online Transaction Processing --- the kind of workload made of many
    small, fast read/write transactions (as opposed to big analytical queries).
\end{terms}
\begin{sayit}
Here's the reassuring part: the database layer doesn't need to change. PostgreSQL is
what Supabase runs today, and what OLK already depends on --- open-source, portable,
proven. Our whole application layer is already free to redeploy anywhere. But it's worth
knowing what China's ecosystem offers if we ever outgrow a single machine: openGauss,
which is actually Huawei's own fork of PostgreSQL, built for distributed, petabyte-scale
workloads. TiDB, from a company called PingCAP, which speaks the same wire protocol as
MySQL but scales horizontally. And OceanBase, from Ant Group, built for financial-grade
transaction processing at extreme concurrency. Because openGauss is a PostgreSQL fork,
that migration path exists --- but we don't need it urgently. We'd only reach for
something like this once a single well-specced machine, like our pilot, stops being
enough --- likely well past a hundred thousand users.
\end{sayit}

\subsection*{Slide 25 --- Supabase self-hosting \timing{1:05}{23:35}}
\onscreen{Three bullets on what Supabase actually is; the ``headline point'' takeaway.}
\begin{terms}
  \item[PostgREST] Software that automatically turns a PostgreSQL database into a
    web API.
  \item[GoTrue] Supabase's authentication service (sign-up, login, sessions).
  \item[Docker] A tool for packaging software (and everything it needs to run) into
    portable containers that run identically on any machine.
  \item[Apache / MIT license] Open-source software licenses that permit free use,
    modification, and redeployment --- including for our own commercial or community
    use.
\end{terms}
\begin{sayit}
This is the headline point of the whole talk, honestly. Supabase itself is just a thin,
open-source layer sitting on top of Postgres --- a piece called PostgREST that turns the
database into an API, GoTrue for authentication, a realtime engine, a storage API, and a
management dashboard called Studio. All open-source, all licensed permissively, all
packaged as Docker containers. We already run this exact stack locally for development.
Running it in production, for real, is the same containers --- just pointed at our own
database, on our own machine, behind our own domain name. We do not need to rewrite Open
Library Kashmir to make it sovereign. We need to redeploy it --- first onto a VPS, then
onto our own hardware.
\end{sayit}

\subsection*{Slide 26 --- Skills matrix \timing{1:05}{24:40}}
\onscreen{Table: skill, why it matters, whether we have it today.}
\begin{terms}
  \item[DBA] Database Administrator --- the role responsible for keeping a database
    backed up, tuned, and running reliably.
  \item[Firewalling] Configuring network rules that control what traffic is allowed
    in and out of our servers.
\end{terms}
\begin{sayit}
None of this runs itself --- here's what our team needs to build. Linux and openEuler
administration, since everything runs on it. Docker and container orchestration ---
we're already good here. PostgreSQL database administration --- backups, tuning,
replication --- we're partway there. Networking and firewalling for shutdown-resilient
connections, security practices like patching and encryption, and basic facilities
knowledge for power and cooling --- those three, honestly, we don't have yet. This is as
much a call for volunteers tonight as it is a call for hardware budget. Every gap in
that ``no'' column is a role somebody in this room could take on.
\end{sayit}

\subsection*{Slide 27 --- Backups \& disaster recovery \timing{1:05}{25:45}}
\onscreen{The 3-2-1 rule adapted for us; encryption takeaway.}
\begin{terms}
  \item[3-2-1 rule] A standard backup principle: 3 copies of data, on 2 different
    media/locations, with 1 copy stored offsite.
  \item[OSS / COS] Alibaba's ``Object Storage Service'' and Tencent's ``Cloud Object
    Storage'' --- cheap, durable cloud storage for files, ideal for encrypted backup
    archives (not for hosting a live database).
\end{terms}
\begin{sayit}
Backups follow a simple, well-tested rule: three-two-one. Three copies of the data ---
the live production copy plus two backups. Two different kinds of media or locations
--- say, a local snapshot plus a separate physical disk or machine. And one copy offsite
--- either a different city in Jammu and Kashmir, or an encrypted bucket on a Chinese
cloud storage service like Alibaba OSS or Tencent COS, used strictly as cold backup,
never as the live host. Everything is encrypted before it ever leaves our building. Used
that way, an offsite Chinese cloud bucket carries none of the sovereignty compromise of
actually hosting live traffic there.
\end{sayit}

\section{The Hope}

\subsection*{Slide 28 --- Divider: The Hope \timing{0:15}{26:00}}
\begin{sayit}
And that brings me to the part I actually want to end on --- the hope.
\end{sayit}

\subsection*{Slide 29 --- The proposal: one machine, real capacity \timing{1:35}{27:35}}
\onscreen{Pilot spec bullets, ``why this scope'', and the headroom calculation row.}
\begin{sayit}
Here's the ask, in concrete terms. A pilot serving ten to twenty thousand users, off a
single, well-specced machine: five terabytes of usable NVMe storage in RAID ten, a
hundred twenty-eight gigabytes of ECC memory, a Kunpeng-class processor, openEuler, and
our own PostgreSQL and self-hosted Supabase stack --- Chinese hardware and hosting from
the virtual machine all the way down to the silicon. This is deliberately modest: small
enough to fund and run with the volunteer skills we can realistically build this year,
but large enough to be a genuinely real platform, not a demo. It proves every single
claim I've made tonight before we ever ask for datacenter-scale money. And the math
backs it up --- twenty thousand users at roughly five megabytes each is only about a
hundred gigabytes of active data. Against five terabytes, that's fifty times headroom.
Years of growth, every book cover, every backup, without touching the hardware again.
\end{sayit}

\subsection*{Slide 30 --- The road from here (closing) \timing{1:35}{29:10}}
\onscreen{Three-phase roadmap diagram; the closing ask; ``Our data. Our servers. Our
region. Let's begin.''}
\begin{terms}
  \item[N+1 redundancy] Having one more of a critical component than the minimum
    needed, so a single failure (of a power supply, a link, a cooling unit) doesn't
    take the system down.
\end{terms}
\begin{sayit}
So here's the road ahead, in three phases. Phase one, this year and next: this exact
pilot, on Chinese VPS or colocation, serving ten to twenty thousand users --- proving the
model works. Phase two, the year after: our own room, with real redundant power and
cooling. Phase three, from twenty-twenty-eight onward: a multi-site data center here in
Kashmir, serving well past a hundred thousand users, as a genuine regional model. What
I'm asking for tonight is simple: a funding commitment for the phase-one parts list,
volunteers against the skills we just went through, and eventually, a room we can call
our own. This is buildable in months, not decades. It starts the moment we agree to
start. Our data. Our servers. Our region. Let's begin. Thank you.
\end{sayit}

\clearpage
\section*{Appendix: Full glossary, A--Z}
\begin{description}[leftmargin=3.2cm,style=sameline,font=\normalfont\bfseries\color{sovDark}]
\item[3-2-1 rule] Backup principle: 3 copies, on 2 media/locations, 1 offsite.
\item[ARM64] A power-efficient processor architecture used by Kunpeng and most mobile
  and modern server chips, as an alternative to x86.
\item[Auto-transfer switch] Hardware that automatically switches power source from
  grid to generator without human intervention.
\item[BOM] Bill of Materials --- an itemised hardware parts list.
\item[Capex] Capital expenditure --- upfront spending on physical assets.
\item[CDN] Content Delivery Network --- cached copies of content served from
  locations closer to the user.
\item[CLOUD Act] US law compelling US companies to hand over data on request
  regardless of where the server is physically located.
\item[Colocation (``colo'')] Renting rack space/power/network in someone else's
  facility while owning the server hardware yourself.
\item[Concurrent users] Users actively using a system at the same moment.
\item[CVM] Cloud Virtual Machine --- Tencent Cloud's rentable virtual server product.
\item[DBA] Database Administrator.
\item[DIMM] Dual In-line Memory Module --- a physical memory stick.
\item[Docker] A tool for packaging software into portable, self-contained containers.
\item[ECC RAM] Error-Correcting Code memory --- detects and fixes single-bit errors.
\item[ECS] Elastic Compute Service --- Alibaba's/Huawei's rentable virtual server
  product.
\item[FM-200 / Novec] Clean-agent fire-suppression gases safe for use around live
  electronics.
\item[Free-air / economizer cooling] Cooling a server room using outside air directly
  instead of mechanical air-conditioning.
\item[Gbps / Mbps] Gigabits/megabits per second --- network speed units.
\item[GbE] Gigabit Ethernet --- wired network speed (e.g.\ 10 GbE = 10 Gbps).
\item[GoTrue] Supabase's authentication microservice.
\item[Hyperscaler] A cloud provider operating at massive global scale.
\item[ICP license] Mainland Chinese filing required only for websites physically
  hosted inside mainland China; avoided by using a Hong Kong/Singapore region.
\item[IOPS] Input/Output Operations Per Second --- how many read/write operations a
  storage device can handle per second.
\item[KPDCL] Kashmir Power Distribution Corporation Limited.
\item[Kunpeng] Huawei's family of ARM-based server processors.
\item[Latency] Round-trip delay for data to travel between two points, in
  milliseconds.
\item[Lighthouse] Tencent Cloud's simplified all-in-one VPS product.
\item[Mirroring] Writing identical data to two drives simultaneously (RAID 1's basis).
\item[N+1 redundancy] Having one spare unit beyond the minimum needed so a single
  failure doesn't cause an outage.
\item[Next.js / React] The open-source web framework and library OLK's front end is
  built with.
\item[NVMe] A storage protocol letting SSDs communicate over high-speed PCIe lanes.
\item[OceanBase] Ant Group's distributed, financial-grade database.
\item[OLTP] Online Transaction Processing --- many small, fast transactions.
\item[openEuler] Huawei's open-source, ARM64-native Linux distribution.
\item[openGauss] Huawei's open-source PostgreSQL fork for distributed workloads.
\item[OSS / COS] Alibaba Object Storage Service / Tencent Cloud Object Storage ---
  cloud file storage, suitable for encrypted backups only.
\item[Parity] Derived data stored alongside real data, allowing reconstruction of a
  failed drive.
\item[Pay-as-you-go] Billing by actual usage rather than a fixed upfront commitment.
\item[PDU] Power Distribution Unit --- rack-mounted power strip.
\item[PostgREST] Software that exposes a PostgreSQL database as a web API.
\item[PostgreSQL] The open-source relational database engine underlying OLK and
  Supabase.
\item[PSU] Power Supply Unit.
\item[RAID] Redundant Array of Independent Disks --- combining drives for redundancy
  and/or speed.
\item[RDIMM] Registered DIMM --- enterprise-grade server memory.
\item[RDSS] Revamped Distribution Sector Scheme --- India's power-sector reform
  program.
\item[SATA SSD] A solid-state drive using the older, slower SATA interface (capped
  $\sim$550 MB/s).
\item[\texttt{shared\_buffers}] A PostgreSQL setting for how much RAM it reserves as
  its own cache.
\item[Striping] Splitting data across multiple drives for parallel, faster
  read/write.
\item[Supabase] Open-source platform bundling Postgres with auth, storage, and an
  API layer.
\item[TaiShan / Inspur / XFusion] Chinese server manufacturers building
  Kunpeng/openEuler-certified rack servers.
\item[TiDB] PingCAP's open-source, MySQL-wire-compatible distributed SQL database.
\item[UPS (online double-conversion)] A battery backup providing zero-delay
  switchover when grid power fails.
\item[VPS] Virtual Private Server --- a rented slice of a shared physical machine.
\item[WAL] Write-Ahead Log --- PostgreSQL's crash-safety mechanism.
\item[YMTC / ZhiTai] China's domestic NAND flash manufacturer and its SSD brand.
\item[2U] A rack server height unit (2 rack units $\approx$ 3.5 inches).
\end{description}

\end{document}
